\documentclass[a4paper,12pt]{article}
\usepackage{docsTemplate}

\setTitle{Specifica Tecnica}
\setAuthors{}
\setVerificators{}
\setApprovation{}
\setVersion{0.0}
\setType{Esterno}
\setDestination{Team Three Technologies \\ & Prof. Tullio Vardanega \\ & Prof. Riccardo Cardin \\ & Sanmarco Informatica SpA}

\begin{document}
    \include{docTitlepage}

    \section*{Registro delle modifiche} {
        \begin{table}[h!]
            \rowcolors{2}{lightgray}{white}
                \begin{tabularx}{\textwidth}{|l|l|l|l|L|L|}
                \hline
                \textbf{Vers.} & \textbf{Data} & \textbf{Autore} & \textbf{Ruolo} & \textbf{Verifica} & \textbf{Oggetto} \\
                \hline
            \end{tabularx}
        \end{table}
    }
    \newpage
    
    \tableofcontents
    \newpage

    \listoffigures
    \newpage
    
    \listoftables
    \newpage

    \section{Introduzione} {
        \subsection{Scopo del documento} {        
            ...
        }

        \subsection{Glossario} {
            Si rimanda al documento Glossario per chiarire i termini nel dominio del progetto che possono risultare ambigui o di difficile comprensione. I vocaboli presenti nel glossario sono stati segnalati seguendo la notazione:
            \begin{center}
                parola\textsubscript{G}
            \end{center}
        }
        
        \subsection{Riferimenti} {
            \subsubsection{Riferimenti normativi} {
                \begin{itemize}
                    \item
                \end{itemize}
            }
            \subsubsection{Riferimenti informativi} {
                \begin{itemize}
                    \item \textit{Glossario v0.5}
                    \begin{itemize}
                        \item \url{https://team-three-technologies.github.io/Documenti/1%20-%20RTB/Documenti%20interni/Glossario_v0_5.pdf}
                        \item[] (Ultimo accesso: 2026/MM/DD)
                    \end{itemize}
                \end{itemize}
            }
        }
    }

    \section{Tecnologie} {

    }

    \section{Architettura} {
        \subsection{Architettura di deployment} {

        }

        \subsection{Architettura logica} {
            
        }

        \subsection{Design pattern} {
            \subsubsection{Dependency injection} {
                
            }
        }
    }

    \section{Stato dei requisiti funzionali} {
        
    }
\end{document}